\chapter{2013 Nobel Prize in Economics}
\textbf{Laureates: }Eugene Fama (USA), Lars Hansen (USA), Robert Shiller (USA)

\section*{Reason for Award}
For their empirical analysis of asset prices

\section{What They Did}
Eugene Fama, Lars Hansen, and Robert Shiller advanced our understanding of asset
prices through complementary approaches that have shaped modern finance. Fama
is known as the "father of modern finance" for his work on efficient markets.
Hansen developed powerful econometric methods for testing asset pricing models.
Shiller documented asset price bubbles and behavioral explanations for price
movements.

Fama developed the efficient market hypothesis, showing that stock prices reflect
all available information and follow random walks. He identified factors like
size, value, and momentum that explain cross-sectional returns. His work on
portfolio diversification and risk management became foundational for investment
practice.

Hansen developed the generalized method of moments (GMM) and stochastic discount
factor framework, providing tools for testing asset pricing models. His 1982
paper introduced GMM, which became one of the most important econometric
methods in economics. The stochastic discount factor approach provides a unified
framework for asset pricing.

His work showed how consumption growth and risk aversion determine asset prices
through the Euler equation. The consumption-based CAPM relates asset returns
to their covariance with consumption growth, providing a theoretical foundation
for risk pricing.

Shiller documented asset price bubbles and excess volatility, challenging the
efficient market hypothesis with evidence that prices move more than
fundamentals justify. He developed measures of stock market valuations and
identified behavioral explanations for price movements.

Shiller's CAPE ratio (cyclically adjusted price-to-earnings ratio) has become
a standard tool for assessing market valuation. His work showed that stock
prices exhibit excess volatility relative to dividends, contradicting the
efficient market hypothesis.

Together, they created the modern framework for empirical asset pricing. Fama's
factor models, Hansen's econometric methods, and Shiller's behavioral insights
combine to form a comprehensive understanding of asset pricing.

\section{Impact on Economic Thinking}
Fama, Hansen, and Shiller transformed finance and macroeconomics. Fama's
efficient market hypothesis became a central paradigm, influencing investment
practice and regulation. His factor models provided tools for portfolio
management and performance evaluation.

The efficient markets thesis reshaped how economists think about information and
price discovery. It influenced debates about active versus passive investing,
with Fama arguing that active management cannot consistently beat the market
after costs.

Hansen's GMM and stochastic discount factor approaches revolutionized econometric
testing, providing flexible methods for evaluating economic models. The Hansen-
Singleton estimator became standard in empirical finance. His work showed how
to test asset pricing models using moment conditions derived from economic
theory.

Shiller's excess volatility findings challenged efficient markets, sparking
debates that enriched financial theory. His behavioral finance research showed
how investor psychology affects prices, leading to the development of
behavioral finance as a field.

Their combined work established asset pricing as an empirical science, with
rigorous testing methods and theoretical foundations. The debate between
efficient markets and behavioral finance continues to drive research, with
evidence supporting both perspectives in different contexts.

The Nobel Prize recognized the complementary nature of their contributions:
Fama's theory, Hansen's methods, and Shiller's empirical challenges.

\section{Modern Perspectives Today}
Asset pricing theory built on their work remains central to finance and economics.
Factor models developed by Fama and French extend his original work, incorporating
profitability and investment dimensions. The five-factor model explains cross-
sectional returns using market, size, value, profitability, and investment factors.

The stochastic discount factor framework underpins modern derivative pricing and
risk management. It provides a unified approach to pricing all assets, from
stocks to options to derivatives.

Shiller's bubble indicators and valuation metrics are widely used by investors
and policymakers. The Shiller P/E ratio has become a standard benchmark for
market valuation, used to assess whether markets are overvalued or undervalued.

Behavioral finance, inspired by Shiller's work, has influenced investment
products like target-date funds and robo-advisors. The 2008 financial crisis
renewed interest in asset pricing anomalies and market efficiency debates.

Experimental and neuroeconomic research tests behavioral hypotheses about
investor decision-making. Climate risk and ESG investing create new asset
pricing challenges that build on their methodological foundations.

Their work continues to influence how economists understand market efficiency,
price formation, and the relationship between risk and return.
